% Class-compatibility preamble retained for a reader excerpt of
%   Dave Witte Morris, "A Cohomological Proof that Real Representations of
%   Semisimple Lie Algebras Have Q-Forms", SIGMA 11 (2015) 034,
%   DOI 10.3842/SIGMA.2015.034, CC BY 4.0.
%
% The source pins the publisher class sigma, which is not part of the delivery
% baseline, so this file supplies the amsart baseline and then copies the author's
% own declaration block mechanically and unchanged:
%   - pinned-original-source.tex lines 11-31: the author's DeclareMathOperator
%     list, the equation numbering and the author's newtheorem declarations;
%   - pinned-original-source.tex lines 35-43, which the author writes immediately
%     after \begin{document} in the publisher layout: the author's own \Res,
%     \ResOne and \Lie macros and \allowdisplaybreaks.
% The publisher header commands of the sigma class (\FirstPageHeading,
% \ArticleName and the like) and the macro \PaperNumber they require are not
% reproduced, because this excerpt carries neither the journal front matter nor
% anything else that uses them.
% No mathematical macro, statement, symbol or assumption is introduced, changed
% or removed by this file.
\documentclass[10pt,a4paper]{amsart}
\usepackage{amssymb,amsmath,amsthm}
\usepackage{enumerate}
\usepackage{hyperref}

\DeclareMathOperator{\Cliff}{Clif\/f}
\DeclareMathOperator{\Gal}{Gal}
\DeclareMathOperator{\Inn}{Inn}
\DeclareMathOperator{\Mat}{Mat}
\DeclareMathOperator{\Out}{Out}
\numberwithin{equation}{section}

\newtheorem{prop}{Proposition}[section]
\newtheorem{cor}[prop]{Corollary}
\newtheorem{lem}[prop]{Lemma}
{\theoremstyle{definition}
\newtheorem{defn}[prop]{Definition}
\newtheorem{notation}[prop]{Notation}
\newtheorem{rem}[prop]{Remark}
\newtheorem{ques}[prop]{Question}
\newtheorem{eg}[prop]{Example}
\newtheorem{erratum}[prop]{Erratum}}

\newcommand{\Res}[1]{\mathop{\mathbf{R}_{#1}}}
\newcommand{\ResOne}[1]{\mathop{\mathbf{R}^{(1)}_{#1}}}
\newcommand{\Lie}[1]{\mathfrak{\lowercase{#1}}}
\allowdisplaybreaks
\usepackage{cleveref}
\crefname{originallemma}{original lemma}{original lemmas}
\Crefname{originallemma}{Original lemma}{Original lemmas}
\makeatletter
\newlabel{SameFixed}{{2.3}{source}{}{sourcepointer-1}{}}
\newlabel{SameFixed@cref}{{[originallemma][0][]2.3}{[1][0][]source}{}{}{}}
\makeatother
\crefname{originallemma}{original lemma}{original lemmas}
\Crefname{originallemma}{Original lemma}{Original lemmas}
\makeatletter
\newlabel{RunivIff}{{2.1}{source}{}{sourcepointer-2}{}}
\newlabel{RunivIff@cref}{{[originallemma][0][]2.1}{[1][0][]source}{}{}{}}
\makeatother
\crefname{originaldefinition}{original definition}{original definitions}
\Crefname{originaldefinition}{Original definition}{Original definitions}
\makeatletter
\newlabel{RunivIff-D}{{4}{source}{}{sourcepointer-3}{}}
\newlabel{RunivIff-D@cref}{{[originaldefinition][0][]4}{[1][0][]source}{}{}{}}
\makeatother
\title{A direct sum of simple Lie algebras over Q is R-universal exactly under four Brauer-theoretic conditions}
\author{Dave Witte Morris}
\date{doi:10.3842/SIGMA.2015.034; CC BY 4.0}
\begin{document}
\maketitle
\noindent\textit{Editorial note.} All delivered mathematical text is the author's own, copied from the pinned publisher source sources/algebra-fl02-a-cohomological-proof-that-real-re/pinned-original-source.tex; no statement or proof step was rewritten, reordered or removed.\par\medskip
\noindent\textit{Editorial note.} Preamble provenance: the source pins the publisher class sigma, which is not part of the delivery baseline (supported wrapper classes are article and amsart), so this card uses sources/algebra-fl02-a-cohomological-proof-that-real-re/retained-preamble.tex. That file supplies the amsart baseline, then copies the author's DeclareMathOperator list, equation numbering and newtheorem declarations from pinned source lines 11-31, and the author's own \textbackslash{}Res, \textbackslash{}ResOne and \textbackslash{}Lie macros from pinned source lines 35-43, which the author writes immediately after \textbackslash{}begin\{document\} in the publisher layout. The publisher front-matter commands of that class and the macro \textbackslash{}PaperNumber they require are not reproduced because this excerpt carries neither the journal front matter nor anything that uses them. The two operator names the author declares that the delivery baseline already supplies (\textbackslash{}GL and \textbackslash{}SL) are left to the baseline, so that no declaration conflicts; their printed meaning is unchanged. No macro was invented and no mathematical symbol was changed.\par\medskip
\noindent\textit{Editorial note.} Bibliography provenance: the counted statement and its proof contain no citation at all, so this card carries no bibliography. The source does have a thebibliography environment (pinned source lines 954-1016) with three entries used elsewhere in the paper; nothing from it is needed here and no entry was reconstructed.\par\medskip
\noindent\textit{Editorial note.} External references are navigation pointers only; every deliberately kept step of the author's argument stays inside the counted spans. Their printed numbers are the author's own theorem numbering under the author's \textbackslash{}newtheorem declarations, computed mechanically by the lane's scanner.\par\medskip
\noindent\textit{Editorial note.} The author writes the proposition label as a LaTeX comment on source line 886 (the line reads percent-backslash-label-SS), so the environment is unlabelled in the author file; the proof refers to the four numbered conditions by their item labels (SS-Gi, SS-star, SS-Di=Dj, SS-Dsplit), all of which lie inside the counted statement span and therefore resolve inside the delivered material.\par\medskip
\noindent\textit{Allowed standard prerequisites (editorial).}\begin{itemize}
\item Brauer equivalence of central simple algebras, denoted by the author with the symbol is-Brauer-equivalent-to, and the classification of real division algebras. Used inside the counted proof (source lines 915, 926-932, 938-941); the author introduces the notation in the first line of the proof. Turns the Brauer-equivalence computations for the tensor products D(Lie G, R)(lambda\_i) tensor D(Lie G, R)(lambda\_j) and for the splits-over-L statements into the isomorphism and splitting conclusions of conditions (3) and (4).
\item Tits algebras D(Lie G, K)(lambda) attached to dominant weights and their behaviour under scalar extension. Used inside the counted proof at source lines 926-932 and 938-941; the author sets this up in the source preceding the counted statement. Supplies the objects whose Brauer classes are being compared, and the identity D(Lie G, K)(lambda\_i + lambda\_j) = D(Lie G, L)(lambda\_i + lambda\_j) over the splitting field L.
\end{itemize}\medskip
\ifdefined\hypertarget\hypertarget{sourcepointer-1}{}\fi
\noindent\textit{Source reference.} Original lemma 2.3: sources/algebra-fl02-a-cohomological-proof-that-real-re/pinned-original-source.tex, source lines 471--471. Navigation pointer only: the counted proof refers to this lemma in one line for condition (2), and the card declares the content of that step as a named standard prerequisite instead of reproducing the lemma.\par\medskip
\ifdefined\hypertarget\hypertarget{sourcepointer-2}{}\fi
\noindent\textit{Source reference.} Original lemma 2.1: sources/algebra-fl02-a-cohomological-proof-that-real-re/pinned-original-source.tex, source lines 416--416. Navigation pointer only: the counted proof invokes this lemma once, in the last paragraph for condition (4), for the equivalence between R-universality and the division-algebra condition; that equivalence is declared as a named standard prerequisite.\par\medskip
\ifdefined\hypertarget\hypertarget{sourcepointer-3}{}\fi
\noindent\textit{Source reference.} Original definition 4: sources/algebra-fl02-a-cohomological-proof-that-real-re/pinned-original-source.tex, source lines 432--432. Navigation pointer only: it is the fourth numbered item of lemma RunivIff, which the counted proof appeals to once in the last paragraph for condition (4); the pointer is here so that the author's own cross-reference number resolves, and the mathematical content it names is declared as a named standard prerequisite.\par\medskip
\begin{prop}
%\label{SS}
A~direct sum $\Lie G = \Lie G_1 \oplus \dots \oplus \Lie G_r$ of simple Lie algebras over~$\mathbb{Q}$ is
$\mathbb{R}$-universal if and only if:
\renewcommand{\labelenumi}{{\rm (\theenumi)}}
\begin{enumerate}\itemsep=0pt
\item
\label{SS-Gi}
each $\Lie G_i$ is $\mathbb{R}$-universal,
\item
\label{SS-star}
either $\Lie G$ is inner, or it becomes inner over an imaginary quadratic extension~$L$ of~$\mathbb{Q}$, and the
$*$-action of $\Gal(\mathbb{C}/\mathbb{R})$ is the same as the $*$-action of $\Gal(L/\mathbb{Q})$.
\item
\label{SS-Di=Dj}
we have $D_{\Lie G_i, \mathbb{Q}}(\lambda_i) \cong D_{\Lie G_j, \mathbb{Q}}(\lambda_j)$ whenever $\lambda_i$
and~$\lambda_j$ are $*$-invariant dominant weights of $G_i$ and~$G_j$, such that $D_{\Lie G_i, \mathbb{Q}}(\lambda_i)
\neq \mathbb{Q}$ and $D_{\Lie G_j, \mathbb{Q}}(\lambda_j) \neq \mathbb{Q}$, and
\item
\label{SS-Dsplit}
either $\Lie G$ is inner, or it becomes inner over an imaginary quadratic extension~$L$ of~$\mathbb{Q}$, such that
$D_{\Lie G_i, \mathbb{Q}}(\lambda_i)$ splits over~$L$, for every dominant weight~$\lambda_i$ of every $\Lie G_i$ that is
inner $($over~$\mathbb{Q})$.
\end{enumerate}
\end{prop}

\begin{proof}
We prove only ($\Rightarrow$), but the argument is reversible.
(We use $\sim$ to denote ``is Brauer equivalent to''.)

(\ref{SS-Gi}) Any representation $\rho_i$ of~$\Lie G_i$ extends to a~representation of $\Lie G$ (that is $0$ on the
other simple factors), so it is clear that $\Lie G_i$ must be $\mathbb{R}$-universal.

(\ref{SS-star}) See Lemma~\ref{SameFixed}.
(This means that all $\Lie G_i$ become inner over the same quadratic extension~$L$.)

(\ref{SS-Di=Dj}) We have
\begin{gather*}
D_{\Lie G, \mathbb{R}}(\lambda_i + \lambda_j) \sim D_{\Lie G, \mathbb{R}}(\lambda_i) \otimes_\mathbb{R} D_{\Lie G,
\mathbb{R}}(\lambda_j) \\
\hphantom{D_{\Lie G, \mathbb{R}}(\lambda_i + \lambda_j)}{}
\cong \bigl(D_{\Lie G, \mathbb{Q}}(\lambda_i) \otimes_\mathbb{Q} \mathbb{R} \bigr)
\otimes_\mathbb{R} \bigl(D_{\Lie G, \mathbb{Q}}(\lambda_i) \otimes_\mathbb{Q} \mathbb{R} \bigr) \cong \mathbb{H}
\otimes_\mathbb{R} \mathbb{H} \sim \mathbb{R}.
\end{gather*}
Since $\Lie G$ is $\mathbb{R}$-universal, this implies $D_{\Lie G, \mathbb{Q}}(\lambda_i + \lambda_j) \sim\mathbb{Q}$,
from which we conclude that $D_{\Lie G, \mathbb{Q}}(\lambda_i) \otimes_\mathbb{Q} D_{\Lie G, \mathbb{Q}}(\lambda_j) \sim
\mathbb{Q}$, so $D_{\Lie G, \mathbb{Q}}(\lambda_i) \cong D_{\Lie G, \mathbb{Q}}(\lambda_j)$ (since these are quaternion
algebras).

(\ref{SS-Dsplit}) Let $\lambda_j$ be a~dominant weight of some (outer) simple factor $G_j$ that is not f\/ixed by the
$*$-action.
Then $\lambda_i + \lambda_j$ is not f\/ixed by the $*$-action, so
\begin{gather*}
D_{\Lie G, \mathbb{Q}}(\lambda_i + \lambda_j) = D_{\Lie G, L}(\lambda_i + \lambda_j) \sim D_{\Lie G_i, L}(\lambda_i)
\otimes_L D_{\Lie G_j, L}(\lambda_j) = D_{\Lie G_i, L}(\lambda_i) \otimes_L L \cong D_{\Lie G_i, L}(\lambda_i).
\end{gather*}
Since $D_{\Lie G, \mathbb{R}}(\lambda_i + \lambda_j) = D_{\Lie G, \mathbb{C}}(\lambda_i + \lambda_j) = \mathbb{C}$, and
$\Lie G$ is $\mathbb{R}$-universal, we conclude from Lemma~\ref{RunivIff}(\ref{RunivIff-D}) that $D_{\Lie G_i, L}(\lambda_i) = L$.
Since $D_{\Lie G, \mathbb{Q}}(\lambda_i) \otimes_\mathbb{Q} L \sim D_{\Lie G, L}(\lambda_i)$, this means that $D_{\Lie
G, \mathbb{Q}}(\lambda_i)$ splits over~$L$.
\end{proof}

\end{document}
